\chapter{Aspects of Geometric Analysis Related to Ricci Flow}
\label{ch:chow-ii-geometric-analysis}

\section{Green's functions}

\begin{definition}[Euclidean Green's function]
\label{def:chow-ii-euclidean-green-function}
\source{chow-et-al-ricci-flow-part-ii:page-0406}
\dcref{appe:1.1}
\group{Green Functions}

For $n\ge3$ set
\[
\Gamma(x,y)=\frac{1}{(n-2)n\omega_n}|x-y|^{2-n},
\]
and for $n=2$ set $\Gamma(x,y)=-(2\pi)^{-1}\log|x-y|$.
\end{definition}

\begin{lemma}[Euclidean Green kernel identities]
\label{lem:chow-ii-euclidean-green-identities}
\source{chow-et-al-ricci-flow-part-ii:page-0406,chow-et-al-ricci-flow-part-ii:page-0407}
\uses{def:chow-ii-euclidean-green-function}
\dcref{appe:1.2}
\group{Green Functions}


For $x\ne y$ and $n\ge3$, $\Delta_x\Gamma=\Delta_y\Gamma=0$, while
\[
|\partial_{x^i}\Gamma|\le(n\omega_n)^{-1}|x-y|^{1-n},
\quad |\partial_{x^i}\partial_{x^j}\Gamma|\le\omega_n^{-1}|x-y|^{-n}.
\]
\end{lemma}

\begin{theorem}[Euclidean fundamental solution identity]
\label{thm:chow-ii-euclidean-green-fundamental}
\source{chow-et-al-ricci-flow-part-ii:page-0407,chow-et-al-ricci-flow-part-ii:page-0408,chow-et-al-ricci-flow-part-ii:page-0409}
\uses{lem:chow-ii-euclidean-green-identities}
\dcref{appe:1.3}
\group{Green Functions}


For $f\in C_c^2(\mathbb R^n)$,
\[
\int_{\mathbb R^n}\Gamma(x,y)\Delta f(y)\,dy=-f(x).
\]
Thus $\Delta_{\mathrm{dist}(y)}\Gamma(x,y)=-\delta_x(y)$ and
$Qf(x)=\int\Gamma(x,y)f(y)\,dy$ is a weak, hence classical, solution of
$\Delta Qf=f$.
\end{theorem}

\begin{theorem}[Green function on a closed manifold]
\label{thm:chow-ii-closed-manifold-green-function}
\source{chow-et-al-ricci-flow-part-ii:page-0409,chow-et-al-ricci-flow-part-ii:page-0410,chow-et-al-ricci-flow-part-ii:page-0411,chow-et-al-ricci-flow-part-ii:page-0412,chow-et-al-ricci-flow-part-ii:page-0413,chow-et-al-ricci-flow-part-ii:page-0414}
\uses{thm:chow-ii-euclidean-green-fundamental}
\dcref{appe:1.4}
\group{Green Functions}


On a closed $n$-manifold, $n\ge3$, a cutoff of
$r^{2-n}/((n-2)\omega_n)$ in geodesic polar coordinates gives a parametrix
$H$.  Iterated convolution with $\Delta H$ improves the singularity by one
power at each step; after at most $n$ steps the remainder is bounded and can
be solved away.  The resulting symmetric kernel $G$ satisfies
\[
\int_M G(x,y)\Delta f(y)\,d\mu(y)=-f(x)+\frac1{\operatorname{Vol}M}\int_M f\,d\mu,
\]
with the usual normalization (and the inverse-Laplacian identity on mean-zero
functions).
\end{theorem}

\begin{lemma}[Convolution singularity estimate]
\label{lem:chow-ii-green-convolution-estimate}
\source{chow-et-al-ricci-flow-part-ii:page-0412,chow-et-al-ricci-flow-part-ii:page-0413,chow-et-al-ricci-flow-part-ii:page-0414}
\dcref{appe:1.5}
\group{Green Functions}

If $|\Psi(x,z)|\le C d(x,z)^{a-n}$ and
$|\Theta(z,y)|\le C d(y,z)^{b-n}$ with $0<a,b<n$, then
\[
\left|\int_M\Psi(x,z)\Theta(z,y)\,d\mu(z)\right|
\le\begin{cases}C'd(x,y)^{a+b-n}&a+b<n,\\
C'(1+|\log d(x,y)|)&a+b=n,\\ C'&a+b>n.
\end{cases}
\]
\end{lemma}

\begin{theorem}[Green asymptotics and estimates]
\label{thm:chow-ii-green-asymptotics}
\source{chow-et-al-ricci-flow-part-ii:page-0414}
\uses{lem:chow-ii-green-convolution-estimate}
\dcref{appe:1.6}
\group{Green Functions}


If $(M^n,g)$ is complete noncompact, $n\ge3$, with nonnegative Ricci
curvature and maximum volume growth, then
\[
\lim_{y\to\infty}n(n-2)\theta_\infty G(x,y)d(x,y)^{n-2}=1.
\]
Moreover, for every $\delta>0$ the two-sided bounds stated in Theorem E.3
hold, with the lower density $\theta_x(\delta d(x,y))$ and asymptotic density
$\theta_\infty$.
\end{theorem}

\begin{theorem}[Hopf boundary point lemmas]
\label{thm:chow-ii-hopf-boundary-point}
\source{chow-et-al-ricci-flow-part-ii:page-0415}
\dcref{appe:1.7}
\group{Maximum Principles}

If a closed domain satisfies the interior sphere condition at $x_0$ and
$u\in C^2$ in its interior obeys $\Delta u+\langle X,\nabla u\rangle\ge0$,
with a strict maximum at $x_0$, then $\partial_\nu u(x_0)>0$.  The same
conclusion holds for $\Delta u+\langle X,\nabla u\rangle+cu\ge0$ when
$c\le0$ and $u(x_0)\ge0$, or when $u(x_0)=0$.
\end{theorem}

\section{Positive and fundamental solutions to the heat equation}

\begin{definition}[Fundamental solution]
\label{def:chow-ii-heat-fundamental-solution}
\source{chow-et-al-ricci-flow-part-ii:page-0416}
\dcref{appe:1.8}
\group{Heat Kernels}

A fundamental solution $h(x,y,t)$ is continuous, $C^2$ in $x$, $C^1$ in $t$,
satisfies $(\partial_t-\Delta_x)h=0$, and converges to $\delta_y$ as
$t\downarrow0$ against compactly supported continuous test functions.
\end{definition}

\begin{lemma}[Duhamel's principle on a closed manifold]
\label{lem:chow-ii-duhamel-closed}
\source{chow-et-al-ricci-flow-part-ii:page-0417}
\dcref{appe:1.9}
\group{Heat Kernels}

For $\square=\partial_t-\Delta$, smooth $A,B$ on a closed manifold satisfy
\[
\int_{t_1}^{t_2}\!\int_M[-(\square A)(w,t-s)B+A(w,t-s)(\square B)]
=\int_M[A(w,t-t_2)B(w,t_2)-A(w,t-t_1)B(w,t_1)].
\]
\end{lemma}

\begin{theorem}[Closed-manifold heat kernel]
\label{thm:chow-ii-closed-heat-kernel}
\source{chow-et-al-ricci-flow-part-ii:page-0417,chow-et-al-ricci-flow-part-ii:page-0418,chow-et-al-ricci-flow-part-ii:page-0419,chow-et-al-ricci-flow-part-ii:page-0420}
\uses{def:chow-ii-heat-fundamental-solution, lem:chow-ii-duhamel-closed}
\dcref{appe:1.10}
\group{Heat Kernels}


On a closed manifold the fundamental solution is unique, symmetric, positive,
and has mass one.  If $\mathscr H_t f(x)=\int_M h(x,y,t)f(y)\,d\mu(y)$, then
\[
\mathscr H_{t_2}\mathscr H_{t_1}=\mathscr H_{t_1+t_2},\qquad
h(x,y,t_1+t_2)=\int_M h(x,z,t_2)h(z,y,t_1)\,d\mu(z).
\]
The IVP $(\partial_t-\Delta)u=F$, $u(\cdot,0)=f$, is represented by
\[
u(x,t)=\int_Mh(x,y,t)f(y)d\mu(y)+\int_0^t\!\int_Mh(x,y,t-s)F(y,s)d\mu(y)ds.
\]
\end{theorem}

\begin{corollary}[Heat kernel sup-norm bound]
\label{cor:chow-ii-heat-kernel-sup-bound}
\source{chow-et-al-ricci-flow-part-ii:page-0420}
\uses{thm:chow-ii-closed-heat-kernel}
\dcref{appe:1.11}
\group{Heat Kernels}


For a heat solution with initial value $f$,
$|u(x,t)|\le\max_yh(x,y,t)\int_M|f|\,d\mu$.
\end{corollary}

\begin{definition}[Dirichlet heat kernel]
\label{def:chow-ii-dirichlet-heat-kernel}
\source{chow-et-al-ricci-flow-part-ii:page-0421}
\dcref{appe:1.12}
\group{Heat Kernels}

On a compact manifold with boundary, a Dirichlet heat kernel is a fundamental
solution in the interior, continuous up to the boundary, and zero when its
first argument lies on the boundary.
\end{definition}

\begin{lemma}[Duhamel's principle with boundary]
\label{lem:chow-ii-duhamel-boundary}
\source{chow-et-al-ricci-flow-part-ii:page-0422}
\dcref{appe:1.13}
\group{Heat Kernels}

For a compact manifold with boundary, the closed-manifold Duhamel identity
acquires the boundary flux term
\[
\int_{t_1}^{t_2}\!\int_{\partial M}\left[-(\partial_\nu A)B+A(\partial_\nu B)\right]d\sigma ds.
\]
\end{lemma}

\begin{theorem}[Minimal positive heat kernel on complete manifolds]
\label{thm:chow-ii-minimal-positive-heat-kernel}
\source{chow-et-al-ricci-flow-part-ii:page-0421,chow-et-al-ricci-flow-part-ii:page-0422,chow-et-al-ricci-flow-part-ii:page-0423,chow-et-al-ricci-flow-part-ii:page-0424,chow-et-al-ricci-flow-part-ii:page-0425}
\uses{def:chow-ii-dirichlet-heat-kernel, lem:chow-ii-duhamel-boundary, thm:chow-ii-hopf-boundary-point}
\dcref{appe:1.14}
\group{Heat Kernels}


On every complete noncompact manifold, an exhaustion by compact regular
domains and their Dirichlet kernels gives an increasing limit $h=\sup_i h_i$.
This is a smooth, positive, symmetric, minimal positive fundamental solution.
If $u$ solves the heat equation with uniformly bounded $L^1$ norm, then on
every compact $K\times[t_0,T)$ all spatial derivatives of $u$ are uniformly
bounded (the local estimate used to pass to the limit).
\end{theorem}

\begin{theorem}[Uniqueness under a Ricci lower bound]
\label{thm:chow-ii-heat-kernel-uniqueness-lower-ricci}
\source{chow-et-al-ricci-flow-part-ii:page-0426}
\dcref{appe:1.15}
\group{Heat Kernels}

If a complete noncompact manifold satisfies $\operatorname{Rc}\ge-k$, its
minimal positive fundamental solution is the unique positive fundamental
solution.
\end{theorem}

\begin{theorem}[Cheeger--Yau heat-kernel comparison]
\label{thm:chow-ii-cheeger-yau-heat-comparison}
\source{chow-et-al-ricci-flow-part-ii:page-0426,chow-et-al-ricci-flow-part-ii:page-0427}
\uses{thm:chow-ii-heat-kernel-uniqueness-lower-ricci}
\dcref{appe:1.16}
\group{Heat Kernels}


If $\operatorname{Rc}\ge(n-1)K$, then
\[
h(x,y,t)\ge H_{n,K}(d(x,y),t),
\]
where $H_{n,K}$ is the radial heat kernel of the simply connected space form
of sectional curvature $K$.
\end{theorem}

\begin{theorem}[Heat kernels on products]
\label{thm:chow-ii-heat-kernel-products}
\source{chow-et-al-ricci-flow-part-ii:page-0427,chow-et-al-ricci-flow-part-ii:page-0428,chow-et-al-ricci-flow-part-ii:page-0429}
\dcref{appe:1.17}
\group{Heat Kernels}

If $h_i$ are fundamental solutions on $(M_i,g_i)$, then
$h_1(x_1,t)h_2(x_2,t)$ is the fundamental solution on the Riemannian product.
The same formula gives the Gaussian fiber factor on a flat rank-$k$ vector
bundle, and periodization gives the heat kernel on a flat quotient such as
$S^1$.
\end{theorem}

\section{Li--Yau and Hamilton estimates}

\begin{theorem}[Li--Yau differential Harnack estimate]
\label{thm:chow-ii-li-yau-estimate}
\source{chow-et-al-ricci-flow-part-ii:page-0430,chow-et-al-ricci-flow-part-ii:page-0431}
\dcref{appe:1.18}
\group{Heat Kernel Estimates}

If $\operatorname{Rc}\ge-kg$ and $u>0$ solves the heat equation, then for
$\alpha>1$,
\[
\partial_t\log u-\alpha^{-1}|\nabla\log u|^2
+\alpha\left(\frac n{2t}+\frac{nk}{2(\alpha-1)}\right)\ge0.
\]
The localized estimate on balls has the additional error
$C\alpha R^{-2}(\alpha^2/(\alpha-1)+\sqrt{k}R)$.
\end{theorem}

\begin{theorem}[Path-integrated Harnack inequality]
\label{thm:chow-ii-path-harnack}
\source{chow-et-al-ricci-flow-part-ii:page-0431}
\uses{thm:chow-ii-li-yau-estimate}
\dcref{appe:1.19}
\group{Heat Kernel Estimates}


For $0<t_1<t_2$,
\[
\frac{u(x_2,t_2)}{u(x_1,t_1)}\ge
\left(\frac{t_2}{t_1}\right)^{-\alpha n/2}
\exp\left[-\frac{\alpha d(x_1,x_2)^2}{4(t_2-t_1)}
-\frac{n\alpha k}{2(\alpha-1)}(t_2-t_1)\right].
\]
For $k=0$ and a fundamental solution, letting $t_1\downarrow0$ gives the
Gaussian lower bound $(4\pi t)^{-n/2}e^{-d(x,y)^2/(4t)}$.
\end{theorem}

\begin{theorem}[Hamilton gradient estimate]
\label{thm:chow-ii-hamilton-gradient-estimate}
\source{chow-et-al-ricci-flow-part-ii:page-0432,chow-et-al-ricci-flow-part-ii:page-0433,chow-et-al-ricci-flow-part-ii:page-0434,chow-et-al-ricci-flow-part-ii:page-0435}
\dcref{appe:1.20}
\group{Heat Kernel Estimates}

If $M$ is closed, $\operatorname{Rc}\ge-k_1g$, and $0<u\le A$ solves the
heat equation, then
\[
t|\nabla\log u|^2\le(1+2k_1t)\log(A/u).
\]
The same maximum-principle argument applies to the stated evolving-metric
heat-type equation with potential, yielding the E.31 estimate.
\end{theorem}

\begin{corollary}[Integral-to-pointwise heat bound]
\label{cor:chow-ii-heat-integral-pointwise}
\source{chow-et-al-ricci-flow-part-ii:page-0434,chow-et-al-ricci-flow-part-ii:page-0435}
\uses{thm:chow-ii-hamilton-gradient-estimate}
\dcref{appe:1.21}
\group{Heat Kernel Estimates}


On a closed manifold and for $t\in(0,T]$,
\[
u(y,t)\le Ct^{-n/2}\int_Mu(x,t)\,d\mu(x).
\]
\end{corollary}

\begin{corollary}[Gradient estimate without a pointwise bound]
\label{cor:chow-ii-heat-gradient-integral}
\source{chow-et-al-ricci-flow-part-ii:page-0436}
\uses{cor:chow-ii-heat-integral-pointwise, thm:chow-ii-hamilton-gradient-estimate}
\dcref{appe:1.22}
\group{Heat Kernel Estimates}


If $u>0$ has unit integral, then
\[
t|\nabla\log u|^2\le C\log\left(\frac{B}{t^{n/2}u}\right)
\]
on $M\times(0,T]$.
\end{corollary}

\begin{theorem}[Laplacian estimate for bounded heat solutions]
\label{thm:chow-ii-heat-laplacian-estimate}
\source{chow-et-al-ricci-flow-part-ii:page-0436,chow-et-al-ricci-flow-part-ii:page-0437}
\dcref{appe:1.23}
\group{Heat Kernel Estimates}

If $M$ is closed, $\operatorname{Rc}\ge-k_1g$, and $0<u\le A$ solves the heat
equation, then
\[
\Delta\log u+2|\nabla\log u|^2
\le\frac{k_1e^{k_1t}}{e^{k_1t}-1}
\left(n+4\log\frac Au\right).
\]
\end{theorem}
